\section{Dispositivos Virtuais}

Todos os dispositivos virtuais implementados são acessados pelos convidados por meio de \gls{mmio}. A maioria deles segue a filosofia e as interfaces padronizadas do \textit{virtio}.
Nesse contexto, a interface C++ que cada dispositivo implementado deve fornecer baseia-se em apenas duas operações fundamentais: \textit{read} e \textit{write}.


\subsection{Fluxo de Mensagens: do Dispositivo Virtual ao Dispositivo Físico}

Quando um \gls{so} convidado envia uma mensagem por meio de um dispositivo virtual, essa mensagem percorre o caminho ilustrado na Figura \ref{fig:virtual_device_to_physical_device}.

\vspace{1cm}

\begin{figure}[h]
	\centering
	\begin{tikzpicture}[
	    node distance=1.4cm, >=stealth,
	    block/.style={rectangle, draw, minimum width=6cm, minimum height=1cm, align=center },
	    decision/.style={diamond, draw, aspect=1, align=center},
	    line/.style={draw, ->, thick}
	]
		\node[block] (vgdevice) {Virtual Device - Guest Side};
		\node[block, dashed, below=of vgdevice] (vmessage) {Virtual Message};
		\node[block, below=of vmessage] (vhdevice) {Virtual Device - Host Side};
		\node[block, below=of vhdevice] (pdevice) {Physical Device};
		\node[block, below=of pdevice, dashed] (pmessage) {Physical Message};

		\node[block, dashed, left=3cm of vgdevice] (hmachine) {Host Machine};
		\node[block, below=of hmachine] (vmachine) {VirtualMachine};

		\draw[line] (vgdevice) -- node[midway, fill=white, inner sep=2pt]{Interrupt} (hmachine);


		\draw[line] (vgdevice) -- (vmessage);
		\draw[line] (vmessage) -- (vhdevice);
		\draw[line] (vhdevice) -- (pdevice);
		\draw[line] (pdevice) -- (pmessage);
		\draw[line] (hmachine) -- node[midway, fill=white, inner sep=2pt]{Notify} (vmachine);
		\draw[line] (vmachine) |- node[midway, fill=white, inner sep=2pt]{Wakeup} (vhdevice);
	\end{tikzpicture}
	\caption{Caminho de Uma Mensagem: do Dispositivo Virtual Para o Dispositivo Físico}
	\label{fig:virtual_device_to_physical_device}
\end{figure}

\subsection{Fluxo de Mensagens: do Dispositivo Físico ao Dispositivo Virtual}

Quando um dispositivo físico precisa enviar uma mensagem a um dispositivo virtual, a mensagem percorre o caminho ilustrado na Figura \ref{fig:physical_device_to_virtual_device}.

\vspace{1cm}

\begin{figure}[h]
	\centering
	\begin{tikzpicture}[
	    node distance=1.4cm, >=stealth,
	    block/.style={rectangle, draw, minimum width=6cm, minimum height=1cm, align=center },
	    decision/.style={diamond, draw, aspect=1, align=center},
	    line/.style={draw, ->, thick}
	]
		\node[block, dashed] (pmessage) {Physical Message};
		\node[block, below=of pmessage] (pdevice) {Physical Device};
		\node[block, below=of pdevice] (vhdevice) {Virtual Device - Host Side};
		\node[block, right=of vhdevice, dashed] (vmessage) {Virtual Message};
		\node[block, below=of vmessage] (vgdevice) {Virtual Device - Guest Side};
		\node[block, below=of vhdevice] (vm) {VirtualMachine};
		\node[block, below=of vm] (vcpu) {VirtualCPU};

		\draw[line] (pmessage) -- (pdevice);
		\draw[line] (pdevice) -- (vhdevice);
		\draw[line] (vhdevice) -- node[midway, fill=white, inner sep=2pt]{Interrupt} (vm);
		\draw[line] (vhdevice) -- (vmessage);
		\draw[line] (vmessage) -- (vgdevice);
		\draw[line] (vm) -- node[midway, fill=white, inner sep=2pt]{Interrupt} (vcpu);
		\draw[line] (vcpu) -| node[midway, fill=white, inner sep=2pt]{Interrupt} (vgdevice);
	\end{tikzpicture}
	\caption{Caminho de Uma Mensagem: do Dispositivo Físico Para o Dispositivo Virtual}
	\label{fig:physical_device_to_virtual_device}
\end{figure}





